\begin{statementbox}
\begin{center}\begin{tikzpicture}[font=\small,>=stealth,line join=round]
\draw[gray!40,rounded corners=3pt] (0,0) rectangle (14.2,8.5);\node at (7.1,7.9) {同样连续对称：两次后是否回到原高度？};
\node at (2.4,6.9) {轴轴：$f=\cos\frac{\pi x}2$};\begin{scope}[shift={(1.2666666666666666,4.348571428571429)},x=0.5666666666666667cm,y=1.0285714285714287cm]\draw[->,gray] (-1,0)--(5,0) node[right]{$x$};\draw[->,gray] (0,-1.7)--(0,1.8) node[above]{$y$};\draw[blue,thick,solid,domain=-0.7:4.8,samples=160] plot(\x,{cos(deg(pi*\x/2))});\foreach \a in {0,2}{\draw[violet,dashed] (\a,-1.4)--(\a,1.4);}\end{scope}\node[align=center,text width=4cm] at (2.4,1.7) {轴 $x=0,2$\\$d=2,\ T=4$};\node at (7.050000000000001,6.9) {点点：$g=\sin\frac{\pi x}2$};\begin{scope}[shift={(5.916666666666667,4.348571428571429)},x=0.5666666666666667cm,y=1.0285714285714287cm]\draw[->,gray] (-1,0)--(5,0) node[right]{$x$};\draw[->,gray] (0,-1.7)--(0,1.8) node[above]{$y$};\draw[teal,thick,solid,domain=-0.7:4.8,samples=160] plot(\x,{sin(deg(pi*\x/2))});\fill[orange!85!black] (0,0) circle(2pt);\fill[orange!85!black] (2,0) circle(2pt);\end{scope}\node[align=center,text width=4cm] at (7.050000000000001,1.7) {中心 $(0,0),(2,0)$\\$d=2,\ T=4$};\node at (11.7,6.9) {轴点：同一条余弦线};\begin{scope}[shift={(10.566666666666666,4.348571428571429)},x=0.5666666666666667cm,y=1.0285714285714287cm]\draw[->,gray] (-1,0)--(5,0) node[right]{$x$};\draw[->,gray] (0,-1.7)--(0,1.8) node[above]{$y$};\draw[blue,thick,solid,domain=-0.7:4.8,samples=160] plot(\x,{cos(deg(pi*\x/2))});\fill[orange!85!black] (1,0) circle(2pt);\draw[violet,dashed] (0,-1.4)--(0,1.4);\end{scope}\node[align=center,text width=4cm] at (11.7,1.7) {轴0，中心 $(1,0)$\\$d=1,\ T=4$};
\node[text width=12.6cm,align=center] at (7.1,.65) {三图同尺度：轴轴、同高点点得到2d；轴点半程翻转高度，需4d。};
\end{tikzpicture}\end{center}


\textbf{\textcolor{CStatement}{先看三个条件}}
\begin{itemize}[leftmargin=*,itemsep=0.25em]
  \item 函数 $f(x)$ 对每个实数 $x$ 都有定义，下面的等式对\textbf{所有实数} $x$ 成立。
  \item 两个对称位置的横坐标为 $a,b$，且 $a\neq b$，记 $d=|a-b|>0$。
  \item 如果有\textbf{两个对称中心}，它们必须处于\textbf{同一高度} $y=c$。
\end{itemize}

\medskip
\textbf{\textcolor{CStatement}{三种情形：算出的都是“一个周期”}}
\end{statementbox}\clearpage\begin{statementbox}\begin{description}[style=nextline,leftmargin=2.4em,itemsep=0.7em]
\item[\textbf{情形一：两条对称轴（轴轴）}] 图像关于 $x=a$ 和 $x=b$ 对称：
\[
\left\{\begin{aligned}f(a+x)&=f(a-x),\\f(b+x)&=f(b-x),\end{aligned}\right.
\qquad \Longrightarrow\qquad \boxed{T=2d}.
\]
\textbf{直观理解：}左右照两次镜子，图像被完整搬到相隔 $2d$ 的位置，样子不变。
\item[\textbf{情形二：两个同高度的对称中心（点点）}] 图像关于 $(a,c)$、$(b,c)$ 对称：
\[
\left\{\begin{aligned}
 f(a+x)+f(a-x)&=2c,\\
 f(b+x)+f(b-x)&=2c,
\end{aligned}\right.
\quad\Longrightarrow\quad \boxed{T=2d}.
\]
\textbf{直观理解：}绕第一个点转半圈，函数值上下翻一次；绕第二个点再转半圈，上下又翻回来。
\item[\textbf{情形三：一条对称轴和一个对称中心（轴点）}] 图像关于 $x=a$ 对称，关于 $(b,c)$ 对称：
\[
\left\{\begin{aligned}
 f(a+x)&=f(a-x),\\
 f(b+x)+f(b-x)&=2c,
\end{aligned}\right.
\quad\Longrightarrow\quad \boxed{T=4d}.
\]
\textbf{直观理解：}前两次对称把图像移了 $2d$，却把函数值变成 $2c-f(x)$；再重复两次，函数值才变回 $f(x)$。
\end{description}

\medskip
\textbf{\textcolor{CStatement}{考试最重要的准确说法}}

$T=2d$ 或 $T=4d$ \textbf{一定是一个正周期，但未必是最小正周期}。
例如最小正周期可能是 $T/3$，仍然可以把 $T$ 叫作“一个周期”。如果题目明确要求\textbf{最小正周期}，必须补充检查。
\end{statementbox}
